% Adopted CH-chain continuation; working spec fragment preserved separately.
\begin{annexitem}{annex:p10.lawvere-ch-chain.cardinal-tools}{thm:p10-lawvere-ch-chain-cardinal-tools}{Bounded maximalization, countable families and absorption}
\begin{proof}
Use the actual base choice and well-order construction of
Theorem~\ref{thm:p10-lawvere-bvm-base}. Its recursion object is the
free pointed endomap object, hence an NNO with zero and successor.
Ordinary recursion with parameters follows by transposing into a function
object; induction follows by applying uniqueness to a successor-closed
subobject containing zero. Zero is not a successor: the recursively
defined truth predicate true at zero and false at each successor separates
them. Successor is monic: recursion on \(\omega\times\omega\), starting
at \((0,0)\) and sending \((a,b)\) to \((s(a),a)\), gives a predecessor
map with \(\operatorname{pred}s=1\). An injective endomap of \([n]\)
is surjective, by induction removing one input and its image; the
zero case is empty. Hence this non-surjective successor cannot occur
on a finite object, and \(\omega\) is not finite.
Finite arithmetic and order are their ordinary NNO recursive operations,
not external integers.

The least union/singleton-closed finite-subset predicate exists by the
power-object intersection construction. Finite enumerations belong to it
by induction on their length. Conversely the predicate of admitting a
finite enumeration contains empty and singletons and is union-closed,
using concatenation; hence it contains that least predicate. In a
classical base remove repeated entries from a finite enumeration in their
order of first occurrence, by finite recursion. This gives a bijection
from some \([n]\). Subsets of \([n]\) are finite by the same finite
recursion. These observations justify all finite-support and finite-fiber
uses below.
The same finite counting induction proves that bijective initial segments
have equal lengths, so finite domains have an actual unique cardinal
map into \(\omega\), not merely chosen enumerations.

Here is the bounded maximalization used twice. For an object \(W\),
let \(\mathcal M\subseteq\mathcal P W\) be a hereditary property containing
empty and preserved by unions of inclusion chains. Well-order \(W\).
At a stage \(w\), take the union of earlier selected subsets and adjoin
\(w\) if the resulting subset satisfies \(\mathcal M\); otherwise leave
that union unchanged. Histories on initial segments are functional graphs
in \(W\times\mathcal P W\), with these equations. Two histories agree
on the common segment: a least disagreement would have the same preceding
union and hence the same prescribed decision. The collection of all
domain/graph pairs is bounded by
\(\mathcal P W\times\mathcal P(W\times\mathcal P W)\).
Its union is therefore a functional history on an initial segment.
If an element is missing, the least missing element can be appended by
the displayed decision; chain-union closure guarantees the preceding
union remains in \(\mathcal M\). This contradicts totality of the union
of all histories. The final selected subset is maximal: an excluded
element already failed when added to a smaller selected subset, and
heredity precludes its addition to the final one. Pairwise disjointness
and pairwise incompatibility satisfy exactly these hypotheses: two
members of a chain union lie together in one of its constituent subsets.
This is an explicit construction, not a use of an unproved Zorn principle.

There is an actual bijection \(\omega^2\to\omega\).
Define \(t_0=0,t_{n+1}=t_n+n+1\) by recursion and send
\((a,b)\) to \(t_{a+b}+b\). Induction on \(k\) shows there is a unique
\(n\) with \(t_n\le k<t_{n+1}\): increase \(n\) just when the next
integer reaches \(t_{n+1}\). Then \(b=k-t_n\le n\) and \(a=n-b\)
give its unique inverse. If a family has countable index object \(K\)
and countable fibers \(A_k\), the object of fiberwise injections into
\(\omega\) is epi over \(K\). Base choice splits it, giving an actual
family of such injections. Combine them with \(K\hookrightarrow\omega\)
and the displayed pairing to embed \(\sum_k A_k\) into \(\omega\).
The sum is the total object of a family map, not an assumed external
infinite coproduct. An image of a countable object is countable because
choice splits its image epi and embeds the image back into that object.
This also proves the assertions for countable unions with overlaps.

For infinite \(X\), well-order \(X\) and recursively choose the least
member not yet used at a finite stage. The state is a finite subobject of
\(X\), an element of \(\mathcal P X\); the remaining object is nonempty
at every finite stage, since otherwise \(X\) would be finite.
Extend the state operation to all of \(\mathcal P X\) by leaving a
state unchanged when its remainder is empty. This defines the total
endomap required for ordinary NNO recursion; the empty-remainder branch
is never reached along finite states of an infinite \(X\).
This constructs a mono \(\omega\to X\) using only NNO recursion.
Let \(W\subseteq\mathcal P X\) consist of subsets bijective with
\(\omega\), and apply bounded maximalization to pairwise-disjoint
subfamilies of \(W\). Its maximal family \(M\) is nonempty. The remainder
of its union is finite, for an infinite remainder would contain a new
copy of \(\omega\) disjoint from every member of \(M\).
Choose one member and add this finite remainder to it.
Every resulting block is still bijective with \(\omega\): prepend the
finite remainder to an enumeration of the chosen block and shift the
enumeration by its finite length. Choice selects a family of these
bijections. The disjoint covering blocks therefore give
\(X\cong M\times\omega\), and hence
\[
       X\times\omega\cong M\times\omega^2\cong M\times\omega\cong X.
 \]
All histories, predicates, graphs and parameter families above are
objects of \(\mathcal S\). The same internal proofs work with arbitrary
parameters; epis in a slice split because their underlying sections
automatically commute with the base projection. No external ordinal
recursion, Replacement or external ZFC model was assumed.
\end{proof}
\end{annexitem}

\begin{annexitem}{annex:p10.lawvere-ch-chain.ccc}{thm:p10-lawvere-ch-chain-ccc}{The finite-support delta-system and finite root coloring}
\begin{proof}
First prove the needed delta-system assertion for an uncountable
subobject \(\mathcal A\) of distinct \(n\)-element subsets of an object
\(D\). A delta-system is a subfamily whose distinct members intersect
in the same fixed root. Induct on \(n\); at zero there is at most one
distinct subset, so no uncountable family.
If some \(d\in D\) lies in uncountably many members, delete \(d\)
from those members. Deletion is injective on that family, and produces
distinct \((n-1)\)-element subsets. Induction gives an uncountable
delta-system there; adding \(d\) back adds it to the common root.

Otherwise each incidence fiber \(\mathcal A_d\) is countable.
Bounded maximalization gives a maximal pairwise-disjoint subfamily
\(\mathcal B\subseteq\mathcal A\). If \(\mathcal B\) were countable,
its union \(U\) would be countable by the proved finite/countable-family
lemma. Maximality says every member of \(\mathcal A\) meets \(U\);
thus the incidence family over \(U\), with countable fibers
\(\mathcal A_d\), maps epimorphically onto \(\mathcal A\).
The countable-family and countable-image lemmas would make
\(\mathcal A\) countable, a contradiction. Hence \(\mathcal B\)
is an uncountable delta-system with empty root. This proves the
assertion by internal natural-number induction, not recursion on
uncountable ordinals.

Now suppose an antichain \(A\subseteq P\) were uncountable.
Its fibers by finite domain cardinal form a family over \(\omega\).
Some fiber \(A_n\) is uncountable, since otherwise their total object
would be countable. Map it to its distinct support family
\(\mathcal A\subseteq\mathcal P D\). A support of size \(n\) admits
exactly \(2^n\) binary functions: order its finite enumeration, code its
bits, and prove the binary enumeration by induction. Thus these fibers
are finite. If \(\mathcal A\) were countable, the family lemma would
make \(A_n\) countable. It is therefore uncountable.

Take an uncountable delta-system of those supports and split the epi from
conditions onto supports, selecting one condition for each support.
Its root is finite. There are only finitely many binary assignments on
that root, so one assignment class is uncountable; otherwise the finite
union would be countable. Any two conditions in this class agree on the
root, and their other domains are disjoint. Their union is functional,
has finite domain, and extends both, contrary to the antichain assumption.
An uncountable object has at least two members because zero and singletons
are countable. This is a contradiction. Thus every antichain is countable.

All branches concern internal predicates on finite-support and incidence
objects. Classicality permits the stated countable/uncountable cases,
and base choice provides the particular families of injections and the
support section. The proof consequently establishes the parameterized
internal ccc, rather than checking only global antichains externally.
\end{proof}
\end{annexitem}

\begin{annexitem}{annex:p10.lawvere-ch-chain.constant-sections}{thm:p10-lawvere-ch-chain-constant-sections}{Constant sheaves, dense decisions and the recursion object}
\begin{proof}
The diagonal of \(\Delta A\) is constant equality. Its complement is
constant inequality, because the base is classical. If two unequal ground
values are given at a condition, no restriction changes them into equal
values. Thus the double-negation closure of this diagonal is itself.
The proved separated-object reflector characterization makes
\(\Delta A\to a\Delta A\) monic, and its unit is dense by the
associated-sheaf construction. Pulling that dense mono back along any
section says exactly that the section is ground-valued on a covering
sieve. Such a sieve is the cofinal set of deciding extensions in the
displayed double-negation coverage. The empty sieve never covers a
condition, since that condition is itself an extension to be tested.
Distinct ground values therefore remain distinct. These statements are
internal and commute with all parameter pullbacks.

We derive the internal site clauses needed below, rather than import
external-small-site theorems into the internal base without a bridge.
For a subpresheaf \(R\hookrightarrow G\), Heyting negation at a
generalized \(\gamma\in G_p\) says that every restriction
\(\gamma|q\), \(q\succeq p\), misses \(R_q\). It is the largest
stable predicate disjoint from \(R\), by its displayed restriction
property and the identity restriction. Applying it twice, and using
classical quantifier duality in \(\mathcal S\), gives
\[
 \gamma\in\overline R_p
 \quad\Longleftrightarrow\quad
 \forall q\succeq p\ \exists r\succeq q,\ \gamma|r\in R_r .
\]
For sheaves, existential image is this closure of the pointwise
presheaf image: the reflector preserves the image factorization,
and closure inside a sheaf is the associated sheaf of the image.
Therefore existential forcing has actual cofinal local witnesses.
For \(R\hookrightarrow G\times A\), the predicate at \(p\) testing
all \(q\succeq p\) and all \(a\in A_q\) is the largest subpresheaf
\(C\hookrightarrow G\) with \(C\times A\subseteq R\).
It is closed when \(R\) is closed: membership locally on a cofinal
cover gives, after pulling that cover to each \(q,a\), local
membership in \(R\), and hence membership by closedness of \(R\).
Thus it is exactly the universal quantifier in sheaves.
This proves the internal all-refinement/cofinal-witness clauses with
base parameters. Full image is consequently exactly the internal
formula \(\forall y\,\exists x\, f(x)=y\). All these are identities
of constructed predicates, not tests on global points.

The geometric adjunction is \(a\Delta\dashv\Gamma i\), from the
already constructed internal Kan adjoints and reflector. For an object
\(Z\) of sheaves with a point and an endomap, apply its right adjoint
to obtain a pointed endomap in \(\mathcal S\); the point transposes
using \(u^*1=1\). The unique base recursion map \(\omega\to u_*Z\)
transposes to a unique map \(u^*\omega\to Z\) preserving zero and
successor, by adjunction naturality. This proves the NNO property, not
just its name. Recursion with parameters follows by the function-object
argument of the preceding proof.

We also justify the stronger source Boolean package for this poset
sheaf topos. It is classical by the established double-negation theorem.
For an internal presheaf \(F\), its free-action counit from the family
of all stage elements onto \(F\) is epi: at a stage each element is the
image of its identity-arrow generator. For a poset, this free-action
presheaf is a subobject of the constant presheaf on the total object
of stage elements, since there is at most one arrow between two stages.
Reflecting gives an actual epi cover \(D_F\to F\), with
\(D_F\hookrightarrow u^*I_F\), where \(I_F\) is that total base object.
The reflected epi remains epi because a left adjoint preserves its
kernel-pair coequalizer and the reflector is left exact.

Code this cover and its kernel pair by the complete Boolean algebra
\(u_*T\), using the subobject/adjunction bijection in
Theorem~\ref{thm:p10-lawvere-bvm-presentation}. That proof used Skolem
only to obtain such covers; the covers here have just been constructed
directly. Effective quotients and pulled-back graphs therefore give the
same actual Boolean-matrix presentations, without assuming Skolem here.
For an epi, the section formula of
Theorem~\ref{thm:p10-lawvere-bvm-realization} applies to these matrices:
well-order its source index, disjointify successive output-truth pieces,
and saturate by the source equality matrix. All its hypotheses, including
internally complete Boolean truth and base choice, have now been supplied.
Thus every epi splits, and the proved source choice equivalence gives
support splitting and \((\Pi)\). This is a source BVM over \(\mathcal S\);
it was not inferred from classicality alone.
\end{proof}
\end{annexitem}

\begin{annexitem}{annex:p10.lawvere-ch-chain.no-epi}{thm:p10-lawvere-ch-chain-no-epi}{Deciding antichains and preservation of the entire epi predicate}
\begin{proof}
Use the internal cofinal-witness/all-refinement clauses and epi
translation just proved. Their external-small-site comparison is
Theorem~\ref{thm:p10-mlm-site-clauses}; it is not an assumption
substituting sets for the present internal base.
Work internally in
\(\mathcal S\), with a condition \(p\) and a staged member \(f\) of
the epi object. All following statements allow an arbitrary base parameter
object; they do not presume that sheaf-global maps or points suffice.
The section \(f\) forces that it is a total surjection from \(u^*X\)
to \(u^*Y\), and restriction preserves that assertion.

For each ground \(x\in X\), the section \(f(x)\) is locally ground.
Let \(D_x\) consist of pairs \((q,y)\) with \(q\succeq p\) forcing
\(f(x)=y\). Its conditions are dense above \(p\); at a fixed \(q\)
there is at most one \(y\), by separation of ground values.
This is a bounded base subobject: restrict the staged \(f\) by its
underlying internal action, evaluate it at the constant-unit image
of \(x\), and pull back equality with the constant-unit image of \(y\).
Its ambient object consists of the base parameter, \(X,P,Y\);
there is no external collection of forcing decisions.
Choose a maximal incompatible subfamily of \(D_x\) by bounded
maximalization. Its condition projection is injective, so the internal
ccc makes it countable. It is predense above \(p\): extend any \(r\)
to a deciding condition; if that extension were incompatible with all
chosen members, it could be added. If any deciding condition forces
value \(y\), some chosen member is compatible with it. At a common
extension the two ground output values are equal. Consequently all
possible values form exactly the image \(S_x\subseteq Y\) of the
chosen antichain's value labels. It is nonempty and countable.
Choice supplies the family of these antichains and their countability
injections uniformly in \(x\), by splitting the corresponding epi
object over \(X\); all families and their predicates are bounded objects.

Surjectivity of \(f\) implies \(\bigcup_xS_x=Y\).
Indeed take a ground \(y\); the existential forcing clause gives an
extension with an input section mapping to \(y\). Refine again until
that input is a ground \(x\). Persistence then puts \(y\) in \(S_x\).
This argument uses two actual cofinal covers, not a current-stage
ground-witness assumption.

If \(X\) is infinite, enumerate each nonempty \(S_x\) by \(\omega\).
Concretely choose an injection \(S_x\hookrightarrow\omega\) and a
default \(s_x\in S_x\); at \(n\) take the unique inverse value if \(n\)
is in the image, and otherwise take \(s_x\). Choice and classicality
make this an actual family of maps \(\omega\to Y\).
Their combined map \(X\times\omega\to Y\) is epi by the preceding
union equality. The proved \(X\times\omega\cong X\) gives an epi
\(X\to Y\), contradicting \(\operatorname{Epi}_{\mathcal S}(X,Y)=0\).
If \(Y=0\), the nonempty deciding-value fiber already contradicts
totality; no default value is assumed in that case.

If \(X\) is finite, enumerate it by \([n]\). Starting at \(p\),
successively extend to decide \(f(x_k)\), for \(k<n\). Finite recursion
on the dense deciding covers gives one extension \(q\) and ground values
\(y_k\) for all those inputs. These define an actual ground map
\(g:X\to Y\). It is not epi by the source assumption. In the
classical base, negating the universal full-image predicate yields a
missing \(y\in Y\); choice may select it with parameters.
But \(q\) still forces \(f\) surjective. A further extension gives a
preimage of \(y\); refine its constant input until it is one of the
ground \(x_k\). Its already decided output is \(y_k\ne y\), a
contradiction. The same argument includes \(n=0\). If \(Y=0\) and
this produces no missing value, \(g\) would have been epi already,
contrary to the original zero-epi-object assumption.

Thus every staged epi member is impossible, in every base context.
Internal Yoneda identifies staged sections with maps from the
sheafified representables; equivalently the underlying internal action
has initial fibers at every condition. Since no empty sieve covers,
the initial sheaf has those empty fibers. Therefore the entire epi
object is zero, not merely its set of global elements.
The source hypothesis \(Y\times\omega\cong Y\) has been retained.
The stronger base choice/cardinal tools here actually make the argument
work without needing that additional restriction; no replacement of
the source's typed lemma by an unproved cardinal slogan occurred.
\end{proof}
\end{annexitem}

\begin{annexitem}{annex:p10.lawvere-ch-chain.generic}{thm:p10-lawvere-ch-chain-generic}{The dense evaluation domain and separation of indices}
\begin{proof}
In presheaves on the actual finite-condition poset, let
\(A\hookrightarrow\Delta(\omega\times I)\) have at \(p\) the
coordinates in \(\operatorname{dom}p\). Its evaluation is
\((n,i)\mapsto p(n,i)\in2\). Extension preserves both the domain and
assigned values, so this is an actual subpresheaf and map.
It is dense: for any \(p,n,i\), leave \(p\) unchanged if the coordinate
is assigned, and otherwise adjoin the assignment \((n,i)\mapsto0\).
The new condition has finite domain and extends \(p\); the same
operation works after every further condition. Hence the membership
sieve is cofinal, in every parameter context.

The reflector inverts this dense mono. Compose its inverse with the
reflected evaluation to obtain
\(u^*\omega\times u^*I\to u^*2\). Product preservation identifies
the reflected constant product with this product of constant sheaves.
Transpose, using the established new number object and
\(u^*2=1+1=T_{\widetilde P}\), to obtain the actual \(\alpha\).
No choice of a total generic ground sequence was made.

For unequal ground \(i,j\), a condition uses finitely many natural
coordinates at those two indices. Their projections are finite subsets
of \(\omega\), so choose a fresh \(n\), for example one greater than
the maximum of the whole finite natural-coordinate projection of \(p\)
(take \(n=1\) for empty projection). Finite recursion constructs that
maximum, independently of enumeration. Adjoin the two assignments
\((n,i)\mapsto0,(n,j)\mapsto1\). The coordinates are distinct and
previously unused, so this is a functional finite extension.
It forces the two evaluated sequences to differ at \(n\).
Such distinguishing extensions are cofinal after every condition.

Here is the kernel argument for arbitrary sheaf sections. Constant
inequality in \(I\times I\) has the dense subpresheaf of distinguishing
extensions just constructed. After reflecting, it covers the complement
of the transported diagonal. On it the kernel pair of \(\alpha\)
is empty, since evaluation differs at the displayed coordinate.
That kernel pair is a closed subobject between sheaves. Its pullback
to the complement diagonal is closed and misses a dense subobject,
so it is zero, by the proved closure/sheaf characterization.
The kernel pair contains the diagonal and has no part in its complement;
classicality makes it precisely that diagonal. Thus \(\alpha\) is mono.
Equivalently, locally ground representatives of arbitrary sections give
the same argument, with all covers refined before comparing them.
This proves actual descent and monicity, not just a pointwise picture
of finite assignments.
\end{proof}
\end{annexitem}

\begin{annexitem}{annex:p10.lawvere-ch-chain.comparisons}{thm:p10-lawvere-ch-chain-comparisons}{All epi predicates in the strict internal comparisons}
\begin{proof}
In any classical power-object base, the singleton predicate map
\(Z\to2^Z\) is monic: membership at the singleton itself separates
unequal elements. No map \(Z\to2^Z\) is epi. For a candidate \(g\),
the predicate \(d(z)=\neg g(z)(z)\) is a member of \(2^Z\);
if \(g(k)=d\), evaluating at \(k\) gives a contradiction.
This works with every parameter and proves the zero epi object,
not merely failure of a particular diagonal enumeration.
Thus \(I=2^C\) is a permissible choice. For \(\omega\), singleton
predicates give \(\omega\hookrightarrow C\), and the same diagonal
argument gives \(\operatorname{Epi}_{\mathcal S}(\omega,C)=0\).
The source's absence of a ground epi also gives this zero-object
formulation: in \(\mathcal S\), a nonzero object has support \(1\),
and support splitting gives it a global member. If its epi predicate
were nonzero it would therefore have a global member, an actual epi.
Thus the given source comparison \(C<I\) supplies precisely the
internally parametrized zero-epi hypothesis used here, not an added
stronger cardinal assumption.
Both \(C\) and the stated \(I\) are infinite, since they contain
copies of \(\omega\). The cardinal-tools theorem therefore supplies
their source-required countable-absorption isomorphisms.
Apply the actual ccc and no-epi theorems to obtain
\(\operatorname{Epi}_{\widetilde P}(u^*\omega,u^*C)=0\).
The transported singleton mono, and the proved NNO identification,
give the left strict comparison.

For the middle-to-right mono, transport base evaluation and transpose
it to obtain \(\kappa:u^*C\to(u^*2)^{u^*\omega}\).
To prove monicity without assuming preservation of this exponential,
take the complement of equality of two ground sequences in the base.
By Boolean quantifier duality it is the image of the predicate that
they differ at some \(n\). Explicitly
\(\neg\forall_n(a(n)=b(n))=\exists_n(a(n)\ne b(n))\);
the identity follows from the two order adjunctions and preservation
of complements by pullback. The geometric inverse image preserves
this image, the diagonal and its complementary sum: it preserves
finite limits and kernel-pair coequalizers. Thus on transported
inequality the images under \(\kappa\) differ at an actual member
of \(u^*\omega\). Its kernel pair has no part on that complement
and is exactly the transported diagonal. Hence \(\kappa\) is mono.
This argument does not commute inverse image with universal quantification
or identify transported ground sequences with all new sequences.

Finally use the proved generic mono \(\alpha:u^*I\hookrightarrow2^\omega\).
The object \(I\) is inhabited: the given mono \(C\to I\) transports
the all-zero sequence. Its inverse image therefore has a global point.
Classicality gives a complementary summand to \(\alpha\).
The copairing of identity on \(u^*I\) and the constant chosen point
on that complement defines an actual retraction
\(r:2^\omega\to u^*I\). It is split epi; no sheaf choice theorem
was needed to make it.
The no-epi lemma with \(X=C,Y=I\) gives
\(\operatorname{Epi}_{\widetilde P}(u^*C,u^*I)=0\).
Any parametrized epi \(u^*C\to2^\omega\), not just \(\kappa\),
would compose with \(r\) to give a parametrized epi to \(u^*I\).
Composition therefore maps its epi predicate into that zero object;
strictness of the initial object makes the original epi predicate zero.
This proves the right strict comparison in the exact source sense.

The constructed \(\widetilde P\) is a nondegenerate source BVM:
the preceding constant-section lemma supplies its stronger Boolean
package, and the empty finite condition ensures its dense site is
nonempty with no empty covering sieve. Its internal cardinal statement
has now been proved with actual maps and all epi predicates.
No model of an independently selected set-theoretic axiom system,
truth-transfer interpretation, or positive consistency direction has
been constructed by these arguments.
\end{proof}
\end{annexitem}
